\documentclass[10pt,a4paper]{amsart}
\usepackage[margin=19mm]{geometry}
\usepackage{amsmath,amssymb,mathtools}
\usepackage[scr=rsfs,cal=euler]{mathalfa}
\usepackage{xfrac}
\usepackage[unicode,hidelinks,bookmarks=false]{hyperref}
\newcommand{\ie}{i.e.\ }
\newcommand{\cf}{cf.\ }
\newcommand{\resp}{respectively\ }
\newcommand{\Subsection}{\subsection}
\providecommand{\nobreakdash}{\nobreak\hbox{-}}
\setlength{\emergencystretch}{2em}
\multlinegap0pt
\usepackage[shortlabels]{enumitem}
\newtheorem{lemma}{Lemma}
\usepackage{cleveref}
\crefname{lemma}{Lemma}{Lemmas}
\title{A free boundary model for invasive and native species under shifting climate in the weak competition case}
\author{Phuoc Vinh Dinh, Phuong Le, Tien Dung Nguyen}
\date{Source: arXiv:2609.14679v1 (CC BY 4.0)}
\begin{document}
\maketitle

\noindent\emph{Source and licence.} A free boundary model for invasive and native species under shifting climate in the weak competition case. Version arXiv:2609.14679v1 (CC BY 4.0), distributed by arXiv under the Creative Commons Attribution 4.0 International licence, http://creativecommons.org/licenses/by/4.0/, as recorded in the arXiv licence metadata for this exact version and shown on the abstract page https://arxiv.org/abs/2609.14679v1. Full licence notice: \texttt{licenses/arxiv-2609.14679-CC-BY-4.0.txt}.\par\smallskip

\noindent\textit{Editorial scope.} Counted unit: the article's lemma lm:mono (source0.tex lines 1334--1336). The article's complete original proof of it is delivered with it (source0.tex lines 1338--1452, one \texttt{proof} environment of 115 lines). No mathematics is written, repaired or re-derived: the statement and the proof are the article's own lines, in the article's own notation, and no retained line is substituted or re-flowed.  Retained uncounted as the source-local context the proof needs: lines 1082--1089; lines 1102--1104; lines 1106--1111; lines 1203--1208; lines 1227--1229; lines 1261--1276.  Nothing was omitted from any retained span, so the package carries no omission record.  No retained line contains a citation, so the wrapper carries no bibliography and no bibliography entry.  No retained line contains a figure environment, an image-inclusion command or drawing code, so no image is delivered and the package declares no asset dependency.  Editorial formatting only, in addition: an AMS \texttt{amsart} preamble that restates the article's own declarations and macros used by the retained lines, namely the environments \texttt{lemma} on separate counters, together with the standard packages those lines need.  The retained environments print in this document as Lemma 1 in order of appearance, whereas the article prints the same items with the numbering of its own full document.  The complete native source of the article as distributed in the arXiv e-print for this exact version is bundled unchanged as \texttt{sources/210346/source0.tex}.
		\begin{equation}\label{eqlm1}
			\begin{cases}
				-c \psi' - d_1 \psi'' = (A(x) - b_1\psi - c_1\varphi)\psi, & -\infty < x < L,\\
				-c \varphi' - d_2 \varphi'' = (a_2 - b_2\psi - c_2\varphi)\varphi, & -\infty < x < +\infty,\\
				\psi(-\infty) = u^*, ~ \psi(x)=0 \text{ for } x\ge L, ~ \psi'(x)<0 \text{ for } x\le L,\\
				\varphi(-\infty) = v^*, ~ \varphi(+\infty)=\frac{a_2}{c_2}, ~ \varphi'(x)>0 \text{ for } x\in \mathbb{R}
			\end{cases}
		\end{equation}

\begin{equation}\label{cooper}
	\Theta_1:=\psi,\qquad \Theta_2:=\frac{a_2}{c_2}-\varphi ,
\end{equation}

\begin{equation}\label{eqcoop}
	\begin{cases}
		d_1\Theta_1''+c\Theta_1'+F_1(x,\Theta)=0, & -\infty<x<L,\\
		d_2\Theta_2''+c\Theta_2'+F_2(\Theta)=0, & -\infty<x<+\infty,
	\end{cases}
\end{equation}

	\begin{equation}\label{s0}
		\Psi_{c_0}(x) \le \psi_L(x) \le u^*,
		\qquad
		v^* \le \varphi_L(x) \le \Phi_{c_0}(x) < \frac{a_2}{c_2}
		\qquad (x\in \mathbb{R}),
	\end{equation}

\begin{lemma}\label{lm:phi_infty}
	$\varphi_L'>0$ on $(L,+\infty)$ and $\varphi_L(+\infty)=\frac{a_2}{c_2}$.
\end{lemma}

\begin{lemma}\label{lm:asym}
	Let $(\psi,\varphi)$ be any solution of \eqref{eqlm1}, and write
	$\psi_L,\varphi_L$ for it. Let $\lambda_1$ be the smallest positive root of
	\begin{equation}\label{charac}
		\Lambda(\lambda):=\big(d_1\lambda^2+c\lambda-b_1u^*\big)\big(d_2\lambda^2+c\lambda-c_2v^*\big)-b_2c_1u^*v^* .
	\end{equation}
	Then $\Lambda$ has exactly two positive roots $0<\lambda_1<\lambda_2$, and
	there are constants $k_1,k_2>0$ (depending on $L$) such that, as
	$x\to-\infty$,
	\begin{equation}\label{asym}
		u^*-\psi_L(x)=\big(k_1+o(1)\big)e^{\lambda_1x},
		\qquad
		\varphi_L(x)-v^*=\big(k_2+o(1)\big)e^{\lambda_1x}.
	\end{equation}
	In particular $\lambda_1$ does not depend on $L$.
\end{lemma}

\begin{lemma}\label{lm:mono}
	$\psi_L'<0$ on $(-\infty,L]$ and $\varphi_L'>0$ on $\mathbb{R}$.
\end{lemma}

\begin{proof}
	For $s>0$ set $\psi^s(x):=\psi_L(x+s)$ and $\varphi^s(x):=\varphi_L(x+s)$.
	Since $A$ is nonincreasing, $A(x+s)\le A(x)$, so on $(-\infty,L-s)$
	\[
	-c(\psi^{s})' - d_1(\psi^{s})'' = \big(A(x+s)-b_1\psi^{s}-c_1\varphi^{s}\big)\psi^{s}
	\le \big(A(x)-b_1\psi^{s}-c_1\varphi^{s}\big)\psi^{s},
	\]
	while $-c(\varphi^{s})'-d_2(\varphi^{s})''=(a_2-b_2\psi^{s}-c_2\varphi^{s})\varphi^{s}$
	on $\mathbb{R}$. In the cooperative variables \eqref{cooper}, this says that
	$\Theta^{s}:=(\psi^{s},\frac{a_2}{c_2}-\varphi^{s})$ is a lower solution of
	\eqref{eqcoop} on $(-\infty,L-s)\times \mathbb{R}$, while $\Theta$ is a solution.

	We use the sliding method. Let
	\[
	S:=\Big\{s>0:\ \psi^{\sigma}\le\psi_L \text{ on }(-\infty,L-\sigma]\ \text{ and }\
	\varphi^{\sigma}\ge\varphi_L \text{ on }\mathbb{R},\ \ \forall\sigma\ge s\Big\}.
	\]

	\emph{Strict inequalities near $-\infty$.} By \cref{lm:asym}, for every
	$\varepsilon\in(0,1)$ there exists $X_\varepsilon<0$ such that
	\begin{equation}\label{eq:two_sided}
		\begin{aligned}
			(1-\varepsilon)k_1e^{\lambda_1y}&\le u^*-\psi_L(y)\le (1+\varepsilon)k_1e^{\lambda_1y},\\
			(1-\varepsilon)k_2e^{\lambda_1y}&\le \varphi_L(y)-v^*\le (1+\varepsilon)k_2e^{\lambda_1y},
		\end{aligned}
		\qquad (y\le X_\varepsilon).
	\end{equation}
	If $\sigma>0$ and $\varepsilon$ satisfy $(1-\varepsilon)e^{\lambda_1\sigma}>1+\varepsilon$,
	then for $x\le X_\varepsilon-\sigma$ both $x$ and $x+\sigma$ lie below
	$X_\varepsilon$, and \eqref{eq:two_sided} gives
	\[
	u^*-\psi^{\sigma}(x)\ \ge\ (1-\varepsilon)k_1e^{\lambda_1\sigma}e^{\lambda_1x}
	\ >\ (1+\varepsilon)k_1e^{\lambda_1x}\ \ge\ u^*-\psi_L(x)
	\]
	and, in the same way, $\varphi^{\sigma}(x)-v^*>\varphi_L(x)-v^*$. Hence
	\begin{equation}\label{eq:strict_minus}
		\psi^{\sigma}<\psi_L \quad\text{and}\quad \varphi^{\sigma}>\varphi_L
		\qquad\text{ on } (-\infty,X_\varepsilon-\sigma]
	\end{equation}
	whenever $(1-\varepsilon)e^{\lambda_1\sigma}>1+\varepsilon$. Note that, given
	any $\sigma_*>0$, this condition holds for all $\sigma\ge\sigma_*$ as soon as
	$\varepsilon$ is small enough.

	\emph{$S\ne\emptyset$.} Take $\varepsilon=\frac12$, write $X:=X_{1/2}$ and let
	$\sigma\ge\sigma_1:=\lambda_1^{-1}\ln5$, so that \eqref{eq:strict_minus} holds
	on $(-\infty,X-\sigma]$. We cover the remaining ranges using that
	$\psi_L<u^*$ on $(-\infty,L]$, that $\varphi_L>v^*$ on $\mathbb{R}$, and that, by
	\eqref{s0} and \cref{lm:phi_infty}, $\varphi_L\le\Phi_{c_0}(L)<\frac{a_2}{c_2}$
	on $(-\infty,L]$ while $\varphi_L$ is increasing on $(L,+\infty)$ with
	$\varphi_L(+\infty)=\frac{a_2}{c_2}$. Fix $Z>L$ with
	$\varphi_L>\Phi_{c_0}(L)$ on $[Z,+\infty)$ and set
	$M_0:=\max_{[X,L]}\psi_L<u^*$ and $m_0:=\min_{[X,Z]}(\varphi_L-v^*)>0$.
	\begin{itemize}
		\item For $x\in[X-\sigma,\,L-\sigma]$ we have $x+\sigma\in[X,L]$, hence
		$\psi^\sigma(x)\le M_0$; since $x\le L-\sigma$ and $\psi_L(-\infty)=u^*>M_0$,
		we get $\psi^\sigma(x)<\psi_L(x)$ for $\sigma$ large. As
		$\psi^\sigma\equiv0$ on $[L-\sigma,+\infty)$, the $\psi$-inequality holds on
		all of $(-\infty,L-\sigma]$.
		\item For $x\in[X-\sigma,\,Z-\sigma]$ we have $x+\sigma\in[X,Z]$, hence
		$\varphi^\sigma(x)-v^*\ge m_0$, while
		$\varphi_L(x)-v^*\le\frac32k_2e^{\lambda_1(Z-\sigma)}<m_0$ for $\sigma$ large.
		\item For $x\in[Z-\sigma,\,L]$ we have $x+\sigma\ge Z$, hence
		$\varphi^\sigma(x)>\Phi_{c_0}(L)\ge\varphi_L(x)$.
		\item For $x\ge L$ we have $\varphi^\sigma(x)=\varphi_L(x+\sigma)>\varphi_L(x)$,
		since $\varphi_L$ is increasing on $(L,+\infty)$.
	\end{itemize}
	Hence $S\supset[s_1,+\infty)$ for $s_1$ large.

	\emph{$\bar s:=\inf S=0$.} By continuity, $\bar s\in S$. Suppose
	$\bar s>0$, and fix $\varepsilon\in(0,1)$ so small that
	$(1-\varepsilon)e^{\lambda_1\bar s/2}>1+\varepsilon$; write
	$X_*:=X_\varepsilon-\bar s$, so that, by \eqref{eq:strict_minus},
	\begin{equation}\label{eq:unif_minus}
		\psi^{\sigma}<\psi_L \quad\text{and}\quad \varphi^{\sigma}>\varphi_L
		\qquad\text{ on } (-\infty,X_*] \text{ for every } \sigma\in[\bar s/2,\bar s].
	\end{equation}
	Set $P:=\psi_L-\psi^{\bar s}\ge0$ on $(-\infty,L-\bar s]$ and
	$Q:=\varphi^{\bar s}-\varphi_L\ge0$ on $\mathbb{R}$. By the computation above and the
	mean value theorem, $(P,Q)$ satisfies a linear cooperative system
	\[
	d_1P''+cP'+\gamma_{11}P+\gamma_{12}Q\le 0,\qquad
	d_2Q''+cQ'+\gamma_{21}P+\gamma_{22}Q\le 0
	\]
	with bounded coefficients and $\gamma_{12},\gamma_{21}\ge0$. Since
	$P(L-\bar s)=\psi_L(L-\bar s)>0$ and, by \eqref{eq:unif_minus},
	$P>0$ and $Q>0$ on $(-\infty,X_*]$, the strong maximum principle for
	cooperative systems yields
	\[
	P>0 \text{ on }(-\infty,L-\bar s],\qquad Q>0 \text{ on } \mathbb{R} .
	\]
	Choose $Z_*>L$ with $\varphi_L>\Phi_{c_0}(L)$ on $[Z_*,+\infty)$. On the
	compact sets $[X_*,L-\bar s]$ and $[X_*,Z_*]$ the continuous functions $P$
	and $Q$ are bounded below by a positive constant; since moreover
	$\psi_L>0$ on $(-\infty,L)$ and $\psi^\sigma(L-\sigma)=0$, a standard
	compactness argument provides $\delta\in(0,\bar s/2)$ such that, for every
	$\sigma\in[\bar s-\delta,\bar s]$,
	\[
	\psi^{\sigma}<\psi_L \ \text{ on }[X_*,L-\sigma]
	\qquad\text{and}\qquad
	\varphi^{\sigma}>\varphi_L \ \text{ on }[X_*,Z_*].
	\]
	Outside these sets the required inequalities hold by \eqref{eq:unif_minus}
	(for $x\le X_*$) and, for $x\ge Z_*$, because
	$\varphi^\sigma(x)=\varphi_L(x+\sigma)>\varphi_L(x)$ by the monotonicity of
	$\varphi_L$ on $(L,+\infty)$. Hence $\bar s-\delta\in S$, contradicting the
	definition of $\bar s$.

	Therefore $\bar s=0$, i.e. $\psi_L(x+s)\le\psi_L(x)$ and
	$\varphi_L(x+s)\ge\varphi_L(x)$ for all $s>0$, so $\psi_L$ is nonincreasing
	and $\varphi_L$ is nondecreasing. Applying the strong maximum principle to
	$\psi_L'$ and $\varphi_L'$, which satisfy linear equations obtained by
	differentiating \eqref{eqlm1} (recall $A$ is Lipschitz and nonincreasing, so
	$A'\le0$ a.e.), we conclude that $\psi_L'<0$ on $(-\infty,L]$ and
	$\varphi_L'>0$ on $\mathbb{R}$.
\end{proof}

\end{document}
